\section{Purpose}

This manual explains how to run and play the current \projectname{} demo in the Unity Editor. The delivered coursework build is a compact single-battle prototype rather than a full campaign game, so the manual focuses on the playable battle loop and the runtime HUD.

\section{Opening the Project}

\begin{enumerate}
    \item Open the Unity project located in \code{coursework/}.
    \item Load \code{Assets/Scenes/SampleScene.unity}.
    \item Press the Unity Editor \texttt{Play} button.
\end{enumerate}

When the scene starts correctly, the player first sees the title screen rather than entering combat immediately.

\section{Starting a Battle}

\begin{enumerate}
    \item Wait for the title screen to finish loading.
    \item Press \texttt{Start Battle}.
    \item The game transitions into the battle state and the battle HUD becomes active.
\end{enumerate}

If the battle has already finished, the result screen allows the player to replay the battle or return to the title screen.

\section{Core Interaction Flow}

The intended interaction order during the player turn is:

\begin{enumerate}
    \item select a player unit on the board;
    \item select a card from the hand;
    \item select a valid target tile or enemy unit if the card requires one;
    \item observe the result and continue until the turn ends.
\end{enumerate}

The player can press \texttt{Clear} to reset the current selection and \texttt{End Turn} to pass control to the enemy phase.

\section{Reading the Battle HUD}

\subsection{Top Bar}

The top bar provides the most important turn-state information:

\begin{itemize}
    \item current turn state;
    \item remaining action budget, shown as \texttt{Energy};
    \item current selection state;
    \item prompt text describing what the player should do next.
\end{itemize}

\subsection{Left Panel}

The left panel provides supporting battle information:

\begin{itemize}
    \item unit summary;
    \item battle log;
    \item recent action history.
\end{itemize}

\subsection{Bottom Hand Panel}

The bottom panel displays the current hand of cards. Each card includes:

\begin{itemize}
    \item card name;
    \item action type;
    \item effect summary such as range, damage, movement, or recovery value.
\end{itemize}

When a card cannot be used in the current state, the card shows disabled feedback directly instead of relying only on the battle log.

\subsection{World-space Unit Feedback}

Each unit displays a world-space health bar above the model. This allows the player to read board state without scanning the text log after every action.

\section{Card Reference}

The current fallback starter deck supports the following actions:

\begin{itemize}
    \item \textbf{Move}: move the selected unit up to 2 tiles.
    \item \textbf{Dash}: move the selected unit up to 3 tiles.
    \item \textbf{Strike}: deal 2 damage to an adjacent enemy.
    \item \textbf{Long Strike}: deal 2 damage to an enemy within 2 tiles.
    \item \textbf{Guard}: reduce the next incoming hit against the selected unit.
    \item \textbf{Push}: deal 1 damage and push the enemy backward if possible.
    \item \textbf{Recover}: restore 2 health to the selected unit.
\end{itemize}

\section{Understanding the Turn Loop}

\begin{enumerate}
    \item At the start of the player turn, the game refreshes the action budget and hand state.
    \item The player performs a limited number of card-driven actions.
    \item The player ends the turn.
    \item Enemy units execute their actions automatically.
    \item The game checks whether the battle should continue or end in victory or defeat.
\end{enumerate}

Some battles may introduce reinforcement pressure during play. This is an intended part of the final tactical slice.

\section{Ending the Battle}

When the battle reaches a win or loss condition, the runtime flow changes to the result screen. The result screen provides:

\begin{itemize}
    \item a visible battle outcome;
    \item replay/start-over behaviour;
    \item a return path to the title screen.
\end{itemize}

\section{Troubleshooting}

\begin{itemize}
    \item If the title screen does not appear, verify that \code{SampleScene.unity} is the loaded scene and that the project compiled successfully.
    \item If a card cannot be played, check the current prompt text, selection state, and the card's disabled feedback.
    \item If board state becomes hard to follow, use the top-bar prompt plus the world-space health bars before relying on the battle log.
\end{itemize}
